\section{Implementation in Practice}
\label{sec:implementation}

One implementation exhibits the grammar against an operational workload distant from constitutional emergency powers.
In an incident-response context, every containment action is constructed as a four-tuple:
the substantive action, a time-to-live, a typed rollback witness, and a quorum predicate.\footnote{The
illustrative workload is endpoint containment.
A quarantine carries a constructive path back to the file's prior location and access-control state;
a process suspension carries the inverse signal that returns the process to its running state;
an egress restriction carries the prior network policy.
Quorum tiers scale with severity, requiring a single operator's authorization for low-severity actions
and concurrent device-level and operator-level signatures within a bounded window for higher-severity ones.
Each admitted action emits a signed receipt that is auditable in canonical form.}
The constructibility of the witness is an empirical feature of the operational regime,
not an a priori feature of the action's category (a point that carries over to the legal domains of Part~\ref{sec:cases}).
The implementation is an existence proof for the grammar's buildability;
it is not an argument that the grammar should apply to the legal regimes the Article examines.
